% ===== ElegantBook customizations for 《深入理解 AI Agent》 =====
% Accent color comes from ElegantBook's cyan theme: \structurecolor = RGB(31,186,190)

% Hungarian document labels. ElegantBook uses its English mode as the base
% layout for Latin-script translations, so override the visible names here.
\AtBeginDocument{%
  \renewcommand{\contentsname}{Tartalomjegyzék}%
  \renewcommand{\chaptername}{Fejezet}%
  \renewcommand{\figurename}{Ábra}%
  \renewcommand{\tablename}{Táblázat}%
  \renewcommand{\listfigurename}{Ábrák jegyzéke}%
  \renewcommand{\listtablename}{Táblázatok jegyzéke}%
}

% ── Text & math fonts: Palatino look, matching the author's papers ──
% The papers' preprint style (arxiv.sty) uses mathpazo (Palatino body +
% matching math) with helvet[scaled=.90] sans. These are the XeLaTeX-native
% equivalents: TeX Gyre Pagella (Palatino) text + TeX Gyre Pagella Math,
% and TeX Gyre Heros (Helvetica clone) sans. Must be set here rather than
% relying on the class: pandoc's template resets the fonts to Latin Modern
% after the class runs, so only preamble-level settings take effect.
% Loaded by filename (like the class's own setup): macOS fontconfig doesn't
% index texmf-dist, so family-name lookup would fail.
\usepackage{unicode-math}
\setmainfont{texgyrepagella}[
  UprightFont = *-regular,
  BoldFont = *-bold,
  ItalicFont = *-italic,
  BoldItalicFont = *-bolditalic,
  Extension = .otf,
  Scale = 1.0]
\setsansfont{texgyreheros}[
  UprightFont = *-regular,
  BoldFont = *-bold,
  ItalicFont = *-italic,
  BoldItalicFont = *-bolditalic,
  Extension = .otf,
  Scale = 0.9]
\setmathfont{texgyrepagella-math}[Extension = .otf]

% ── Unicode glyph fallback (the text/mono fonts lack these) ──────
\usepackage{newunicodechar}
% macOS has Arial Unicode MS / Heiti SC; on Linux those Apple fonts are absent,
% so fall back to broad-coverage libre fonts to keep the build working.
\IfFontExistsTF{Arial Unicode MS}{\newfontfamily\fallbackfont{Arial Unicode MS}}{\newfontfamily\fallbackfont{DejaVu Sans}}
\IfFontExistsTF{Heiti SC}{\newunicodechar{★}{{\fontspec{Heiti SC}★}}}{\newunicodechar{★}{{\fontspec{Noto Sans CJK SC}★}}}
\newunicodechar{✓}{{\fallbackfont ✓}}
\newunicodechar{✗}{{\fallbackfont ✗}}
\newunicodechar{₀}{{\fallbackfont ₀}}
\newunicodechar{₁}{{\fallbackfont ₁}}
\newunicodechar{₂}{{\fallbackfont ₂}}
\newunicodechar{₃}{{\fallbackfont ₃}}
\newunicodechar{₄}{{\fallbackfont ₄}}
\newunicodechar{₅}{{\fallbackfont ₅}}
\newunicodechar{₆}{{\fallbackfont ₆}}
\newunicodechar{₇}{{\fallbackfont ₇}}
\newunicodechar{₈}{{\fallbackfont ₈}}
\newunicodechar{₉}{{\fallbackfont ₉}}
\newunicodechar{ⓐ}{{\fallbackfont ⓐ}}
\newunicodechar{ⓑ}{{\fallbackfont ⓑ}}
\newunicodechar{ⓒ}{{\fallbackfont ⓒ}}
\newunicodechar{ⓓ}{{\fallbackfont ⓓ}}
\newunicodechar{ⓔ}{{\fallbackfont ⓔ}}
% Math / symbols absent from the Latin text font
\newunicodechar{≈}{{\fallbackfont ≈}}
\newunicodechar{♠}{{\fallbackfont ♠}}
\newunicodechar{♣}{{\fallbackfont ♣}}
\newunicodechar{♥}{{\fallbackfont ♥}}
\newunicodechar{♦}{{\fallbackfont ♦}}
% Greek letters used inline in body text (Latin Modern lacks them here)
\newunicodechar{α}{{\fallbackfont α}}
\newunicodechar{β}{{\fallbackfont β}}
\newunicodechar{γ}{{\fallbackfont γ}}
\newunicodechar{δ}{{\fallbackfont δ}}
\newunicodechar{ε}{{\fallbackfont ε}}
\newunicodechar{η}{{\fallbackfont η}}
\newunicodechar{θ}{{\fallbackfont θ}}
\newunicodechar{λ}{{\fallbackfont λ}}
\newunicodechar{μ}{{\fallbackfont μ}}
\newunicodechar{π}{{\fallbackfont π}}
\newunicodechar{ρ}{{\fallbackfont ρ}}
\newunicodechar{σ}{{\fallbackfont σ}}
\newunicodechar{τ}{{\fallbackfont τ}}
\newunicodechar{φ}{{\fallbackfont φ}}
\newunicodechar{ω}{{\fallbackfont ω}}

% ── Residual CJK glyphs ──────────────────────────────────────────
% Chinese terminology and tokenization examples occur throughout the Hungarian
% text. Use xeCJK for complete fallback coverage instead of enumerating a
% brittle list of individual code points.
\usepackage{xeCJK}
\IfFontExistsTF{Heiti SC}{%
  \setCJKmainfont{Heiti SC}%
  \setCJKmonofont{Heiti SC}%
}{%
  \setCJKmainfont{Noto Sans CJK SC}%
  \setCJKmonofont{Noto Sans Mono CJK SC}%
}

% ── Code/monospace font (covers box-drawing glyphs used in code) ──
% Menlo is a macOS font; fall back to DejaVu Sans Mono on Linux.
\IfFontExistsTF{Menlo}{\setmonofont[Scale=0.86]{Menlo}}{\setmonofont[Scale=0.86]{DejaVu Sans Mono}}

% ── tcolorbox (ElegantBook already loads it; just add libraries) ──
\tcbuselibrary{skins,breakable}

% ── TikZ (for the full-bleed image cover in cover.tex) ───────────
\usepackage{tikz}
\usetikzlibrary{fadings}
% Vertical fade: invisible at the top, solid at the bottom — lets the cover
% art melt smoothly into the navy title band.
\tikzfading[name=titlefade, top color=transparent!100, bottom color=transparent!0]

% Theme accent color: override ElegantBook's bright cyan with a deep navy
% (classic, authoritative technical-book tone). Drives headings, rules, links,
% code chrome and the cover. Tweak this one line to restyle the whole book.
\definecolor{structurecolor}{RGB}{30,58,107}

% Accent shade for code/figure chrome (follows the theme color above)
\colorlet{accent}{structurecolor}

% Internal cross-reference links (图N-M / 第N章) — clickable + subtle teal,
% independent of the document class's global linkcolor. Used by crossref.lua.
\newcommand{\crossreflink}[2]{\hyperref[#1]{\textcolor{structurecolor}{#2}}}
\colorlet{accentbg}{structurecolor!6}
\colorlet{accentrule}{structurecolor!35}

% ── Wrap long lines in code blocks ───────────────────────────────
% pandoc emits fancyvrb's Verbatim (via `Highlighting` for ```lang blocks
% and plain `verbatim` for un-tagged blocks); neither wraps by default, so
% long lines overflow the page. fvextra adds breaklines; re-define both
% environments to wrap, breaking mid-token when a chunk has no spaces (URLs,
% long identifiers, CJK comments). A teal ↳ marks each continued line.
\usepackage{fvextra}
\fvset{breaklines=true,breakanywhere=true,%
  breaksymbolleft={\footnotesize\textcolor{accentrule}{$\hookrightarrow$}\,}}
\DefineVerbatimEnvironment{Highlighting}{Verbatim}{commandchars=\\\{\},breaklines,breakanywhere}

% ── Code blocks: subtle O'Reilly-style framed box ────────────────
% One shared box style for ALL code blocks so highlighted (```lang → Shaded)
% and plain (``` → verbatim) blocks look identical. The plain verbatim nests
% a fvextra Verbatim inside the box; \VerbatimEnvironment makes it end at
% \end{verbatim} (the outer env name) rather than \end{Verbatim}.
\newtcolorbox{codebox}{%
  breakable, enhanced,
  colback=black!3, colframe=accentrule,
  boxrule=0.5pt, leftrule=2.5pt,
  left=0.7em, right=0.6em, top=0.4em, bottom=0.4em,
  arc=0.6mm,
  before skip=0.7em, after skip=0.7em
}
\renewenvironment{Shaded}{\begin{codebox}}{\end{codebox}}
\renewenvironment{verbatim}%
  {\VerbatimEnvironment\begin{codebox}\begin{Verbatim}}%
  {\end{Verbatim}\end{codebox}}

% ── Experiment & question callouts (used by experiment_box.lua) ───
\newtcolorbox{experimentbox}{
  breakable, enhanced,
  colback=accentbg, colframe=accent,
  boxrule=0.6pt, left=0.8em, right=0.8em, top=0.6em, bottom=0.6em,
  arc=1mm, before skip=0.9em, after skip=0.9em
}

\newtcolorbox{questionbox}{
  breakable, enhanced,
  colback=accentbg, colframe=accent,
  boxrule=0.8pt, left=0.9em, right=0.9em, top=0.9em, bottom=0.7em,
  arc=1.2mm, before skip=0.9em, after skip=0.9em,
  title={\sffamily\bfseries Elgondolkodtató kérdések},
  fonttitle=\normalsize\color{white},
  coltitle=white,
  attach boxed title to top left={xshift=1em, yshift=-\tcboxedtitleheight/2},
  boxed title style={colback=accent, colframe=accent, arc=0.8mm, outer arc=0.8mm}
}

% ── Figures: force [H] so they can sit inside callout/quote boxes ─
\usepackage{float}
\let\origfigure\figure
\let\origendfigure\endfigure
\renewenvironment{figure}[1][htbp]{\origfigure[H]}{\origendfigure}
\setkeys{Gin}{width=\linewidth,height=0.38\textheight,keepaspectratio}

% ── Figure captions: manual numbering in text, so suppress label ──
\usepackage{caption}
% Sans, bold, centered. Figure numbers live in the caption text itself
% (crossref.lua), so the automatic "Figure N:" label is suppressed.
\captionsetup{font={normalsize,sf,bf}, labelformat=empty, skip=4pt, justification=centering}

% ── Blockquotes: teal left bar (experiments render as quotes) ────
\renewenvironment{quote}{%
  \begin{tcolorbox}[
    blanker, breakable,
    borderline west={2.5pt}{0pt}{accent},
    colback=accent!4,
    left=1em, right=0.6em, top=0.5em, bottom=0.5em,
    before skip=0.6em, after skip=0.6em
  ]%
}{%
  \end{tcolorbox}%
}

% ── Inline code: keep it tight ───────────────────────────────────
\setlength{\parindent}{2em}

% ElegantBook empties the plain page style used on chapter-opening pages.
% Keep those pages numbered, matching the centered footer on regular pages.
\fancypagestyle{plain}{%
  \fancyhf{}%
  \fancyfoot[C]{\color{structurecolor}\small\thepage}%
  \renewcommand{\headrulewidth}{0pt}%
  \renewcommand{\footrulewidth}{0pt}%
}
